\documentclass[11pt,leqno]{amsart}
\usepackage[letterpaper,margin=1in]{geometry}
\usepackage[T1]{fontenc}
\usepackage{lmodern}
\usepackage{amsmath,amssymb,mathtools}
\usepackage{microtype}
\usepackage{enumitem}
\usepackage[hidelinks]{hyperref}
\hypersetup{
  pdftitle={A locally sparse graph with superlinear Ramsey growth},
  pdfauthor={Will Howe, Grant Huffman, Jensen Kohlmeyer, Liam Kruer},
  pdfsubject={A counterexample to Erdos problem 566},
  pdfkeywords={Ramsey theory, locally sparse graphs, entropy, Hamming cube}
}
\numberwithin{equation}{section}
\newtheorem{theorem}{Theorem}[section]
\newtheorem{proposition}[theorem]{Proposition}
\newtheorem{lemma}[theorem]{Lemma}
\newtheorem{corollary}[theorem]{Corollary}
\theoremstyle{definition}
\newtheorem{definition}[theorem]{Definition}
\theoremstyle{remark}
\newtheorem{remark}[theorem]{Remark}
\DeclareMathOperator{\Var}{Var}
\DeclareMathOperator{\supp}{supp}
\newcommand{\E}{\mathbb E}
\newcommand{\Prob}{\mathbb P}
\newcommand{\F}{\mathbb F}
\newcommand{\one}{\mathbf 1}
\newcommand{\HH}{\mathrm H}
\newcommand{\dist}{\mathrm d}
\newcommand{\eps}{\varepsilon}
\newcommand{\set}[1]{\{#1\}}
\newcommand{\abs}[1]{\lvert#1\rvert}
\newcommand{\norm}[1]{\lVert#1\rVert}
\newcommand{\floor}[1]{\lfloor#1\rfloor}
\newcommand{\ceil}[1]{\lceil#1\rceil}
\newcommand{\bin}[2]{\binom{#1}{#2}}
\newcommand{\R}{\mathbb R}
\setlist[enumerate]{leftmargin=2em,itemsep=2pt,topsep=5pt}
\setlength{\emergencystretch}{1.5em}
\allowdisplaybreaks[2]

\title[Locally sparse graphs and Ramsey growth]{A locally sparse graph with superlinear Ramsey growth}
\author[Howe, Huffman, Kohlmeyer, and Kruer]{Will Howe, Grant Huffman, Jensen Kohlmeyer, and Liam Kruer}
\date{November 11, 2026}
\subjclass[2020]{05C55, 05D40}
\keywords{Ramsey numbers, locally sparse graphs, entropy, Hamming cube, Krawtchouk polynomials}
\thanks{Electronic addresses are \href{mailto:howew@purdue.edu}{\nolinkurl{howew@purdue.edu}} (W. Howe),
\href{mailto:huffma97@purdue.edu}{\nolinkurl{huffma97@purdue.edu}} (G. Huffman),
\href{mailto:jenwin@purdue.edu}{\nolinkurl{jenwin@purdue.edu}} (J. Kohlmeyer), and
\href{mailto:kruerl@purdue.edu}{\nolinkurl{kruerl@purdue.edu}} (L. Kruer).}

\begin{document}
\begin{abstract}
We construct a fixed finite bipartite graph $G$ such that every set of at least two vertices spans at most twice its cardinality minus three edges, while the ratio $R(G,H)/e(H)$ is unbounded as $H$ ranges over graphs without isolated vertices. For all sufficiently large $m$,
\[
  R\bigl(G,K_{2^{570m}}\bigr)>2^{1144m-1}.
\]
This gives a negative answer to Erd\H{o}s problem 566. The construction uses a layered two-parent graph and a random subset of a Hamming cube. A conditional-entropy inequality excludes the fixed graph, and a bound on the negative Krawtchouk spectrum, followed by an independent-set encoding, excludes large independent sets. We give the complete argument, including the finite scalar estimates.
\end{abstract}
\maketitle

\section{Introduction}

For finite simple graphs $G$ and $H$, let $R(G,H)$ be the least integer $n$ such that every red-blue coloring of the edges of $K_n$ contains a red copy of $G$ or a blue copy of $H$. Copies are subgraphs. Write $e(H)$ for the number of edges of $H$.

Erd\H{o}s, Faudree, Rousseau, and Schelp \cite{EFRS} call a graph $G$ \emph{Ramsey size-linear} if there is a constant $C_G$ such that
\begin{equation}\label{eq:size-linear}
  R(G,H)\le C_G e(H)
\end{equation}
for every finite simple graph $H$ without isolated vertices. They asked whether the local condition
\begin{equation}\label{eq:local-condition}
  e(G[W])\le 2\abs{W}-3
  \qquad\text{for every }W\subseteq V(G)\text{ with }\abs{W}\ge2
\end{equation}
implies \eqref{eq:size-linear}. This question is recorded as Erd\H{o}s problem 566 \cite{Bloom}. The formulation using all subgraphs is equivalent to \eqref{eq:local-condition}, because each subgraph is contained in the induced subgraph on its vertex set.

\begin{theorem}\label{thm:main}
There is a fixed finite bipartite, $2$-degenerate graph $G$ satisfying \eqref{eq:local-condition} such that, for every sufficiently large positive integer $m$,
\begin{equation}\label{eq:main-bound}
 R\bigl(G,K_{2^{570m}}\bigr)>2^{1144m-1}.
\end{equation}
In particular,
\begin{equation}\label{eq:main-ratio}
 \frac{R\bigl(G,K_{2^{570m}}\bigr)}{e\bigl(K_{2^{570m}}\bigr)}>2^{4m}.
\end{equation}
\end{theorem}

The graph $G$ has finitely many layers. Each vertex after the first layer represents a pair of vertices in the preceding layer and is joined to those two vertices. The layered construction, the entropy-potential approach, and the use of profiles to count child arrays originate in Chapter 10 of OpenAI \cite{OpenAI}. We use a different entropy parameter and apply the construction inside a single randomly sampled cube. The independent-set encoding used later is the Kleitman--Winston method, for which \cite[Section 3]{BTW} gives an account. The specific estimates required here are proved below.

We prove that one random vertex set simultaneously supports a $G$-free short-distance graph and has small independence number. Coloring its complement blue then yields \eqref{eq:main-bound}.

\section{Entropy estimates}

All entropies and logarithms denoted by $\log_2$ use base two. Natural logarithms are denoted by $\ln$. Set
\begin{equation}\label{eq:constants}
 \tau=\frac{43}{100},\qquad
 \beta=\frac{4857}{5000},\qquad
 w=\frac9{16},\qquad
 \sigma=-\log_2 w.
\end{equation}
The binary entropy function is
\[
 h(x)=-x\log_2 x-(1-x)\log_2(1-x),\qquad 0\le x\le1,
\]
with the endpoint values defined by continuity. We shall use
\begin{equation}\label{eq:scalar-bounds}
\begin{aligned}
 .94268&<h(.36)<.94270,\\
 .98581&<h(.43)<.98582,\\
 .04464&<h(.0049)<.04465.
\end{aligned}
\end{equation}
These are strict rational bounds. Appendix~\ref{app:scalars} proves them from a logarithmic series with an explicit remainder.

\subsection{Elementary entropy facts}

For a finite random variable $X$ with probabilities $p_x$, write
\[
 \HH(X)=-\sum_x p_x\log_2 p_x.
\]
Conditional entropy is the average of the entropies of the conditional laws. The following elementary estimates also fix the conventions used at zero probabilities.

\begin{lemma}\label{lem:elementary-entropy}
For $a,b\in[0,1]$,
\begin{equation}\label{eq:continuity}
 \abs{h(a)-h(b)}\le h(\abs{a-b}).
\end{equation}
Moreover,
\begin{equation}\label{eq:pinsker-binary}
 h\left(\frac{1+u}{2}\right)\le1-\frac{u^2}{2\ln2}
 \qquad(-1\le u\le1).
\end{equation}
For integers $n\ge1$ and $0\le j\le n$,
\begin{equation}\label{eq:binomial-entropy}
 \frac{2^{n h(j/n)}}{n+1}\le\binom nj\le 2^{n h(j/n)}.
\end{equation}
\end{lemma}

\begin{proof}
Couple Bernoulli variables $A$ and $B$ with means $a$ and $b$ so that $\Prob(A\ne B)=\abs{a-b}$. With $D=\one_{A\ne B}$, each of $A$ and $B$ is determined by the other together with $D$. The entropy chain rule gives $\HH(A)\le\HH(B)+\HH(D)$ and the reverse inequality with $A$ and $B$ interchanged. This proves \eqref{eq:continuity}.

For $0<x<1$,
\[
 h''(x)=-\frac{1}{x(1-x)\ln2}\le-\frac4{\ln2}.
\]
Integrating twice from $x=1/2$, where $h(1/2)=1$ and $h'(1/2)=0$, proves \eqref{eq:pinsker-binary}. Continuity supplies the endpoints.

For $0<j<n$, put $q=j/n$. In the binomial distribution with parameters $n,q$, the term at $j$ is maximal. Indeed, the ratio of the term at $i+1$ to the term at $i$ is
\[
 \frac{n-i}{i+1}\frac{j}{n-j},
\]
which is greater than one for $i<j$ and less than one for $i\ge j$. The maximal term is between $1/(n+1)$ and $1$. Substituting
\[
 \binom nj q^j(1-q)^{n-j}=\binom nj\,2^{-n h(q)}
\]
gives \eqref{eq:binomial-entropy}. For $j=0,n$, the inequalities follow directly.
\end{proof}

\subsection{A two-parent inequality}

Define
\begin{equation}\label{eq:kappa-theta}
 \kappa=\frac12+\frac12\log_2(1+w)+\frac14\log_2 w,
 \qquad \theta=\frac{1+h(9/25)}2.
\end{equation}

\begin{lemma}\label{lem:two-parent}
Let $X,Y$ be independent Bernoulli variables with the same mean $q$. Let $Z$ be a binary random variable with arbitrary conditional law given $(X,Y)$, and write $v=\E Z$ and
\[
 d=\frac{\Prob(X\ne Z)+\Prob(Y\ne Z)}2.
\]
Then
\begin{equation}\label{eq:two-parent}
 \HH(Z\mid X,Y)\le
 \kappa+\sigma d+\frac{h(v)-h(q)}2.
\end{equation}
Furthermore,
\begin{equation}\label{eq:theta-identity}
 \kappa+\sigma\tau=\theta<.97135<\beta.
\end{equation}
\end{lemma}

\begin{proof}
We first take $0<q,v<1$. Put $\pi_z=\Prob(Z=z)$ and
\[
 c(x,y,z)=\frac{\one_{x\ne z}+\one_{y\ne z}}2.
\]
For a fixed parent pair $(x,y)$, apply the entropy bound associated with the positive weights
\[
 t_z=\sqrt{\pi_z}\,w^{c(x,y,z)}.
\]
For any probability vector $(p_z)$ and positive weights $(t_z)$, this bound is
\[
 -\sum_zp_z\log_2p_z
 \le\log_2\Bigl(\sum_zt_z\Bigr)-\sum_zp_z\log_2t_z.
\]
To verify it, normalize $t_z$ to a probability vector and use $\ln t\le t-1$ to obtain nonnegativity of relative entropy. Averaging over $(X,Y)$ therefore gives
\begin{equation}\label{eq:gibbs}
 \HH(Z\mid X,Y)
 \le \sigma d+\frac12h(v)+\E\log_2\Bigl(\sum_z t_z\Bigr).
\end{equation}

Let $u=2q-1$, $b_0=\sqrt{1-v}$, and $b_1=\sqrt v$. The probabilities of parent types $00$, $11$, and the two mixed types together are
\[
 a_0=\frac{(1-u)^2}4,\qquad
 a_1=\frac{(1+u)^2}4,\qquad
 a_* =\frac{1-u^2}2.
\]
It suffices by \eqref{eq:gibbs} to show that
\begin{align*}
 F={}&\frac12h(q)
 +a_0\log_2(b_0+w b_1)
 +a_1\log_2(w b_0+b_1)\\
 &\hspace{35mm}+a_*\log_2\bigl(\sqrt w(b_0+b_1)\bigr)
 \le\kappa.
\end{align*}
Set
\[
 M_0=\frac{1+w}{\sqrt2},\qquad
 M_* =\sqrt{2w},\qquad
 r=\frac{1-w}{1+w}=\frac7{25}.
\]
Apply $\ln z\le\ln M+z/M-1$ to the first two logarithms at $M=M_0$ and the third at $M=M_*$. The weighted logarithmic constants sum to
\[
 \frac12\ln(1+w)+\frac14\ln w
 +\frac{u^2}2\ln\frac{M_0}{M_*}.
\]
The weighted linear terms sum to
\[
 \frac{(1-ru)b_0+(1+ru)b_1}{\sqrt2}-1
 \le\sqrt{1+r^2u^2}-1\le\frac{r^2u^2}2.
\]
Here the first inequality is Cauchy--Schwarz and uses $b_0^2+b_1^2=1$. Also $M_0/M_*=(1-r^2)^{-1/2}$. Combining these observations with \eqref{eq:pinsker-binary} gives
\begin{equation}\label{eq:entropy-bracket}
 (F-\kappa)\ln2
 \le u^2\left(-\frac14-\frac14\ln(1-r^2)+\frac{r^2}2\right).
\end{equation}
Since $-\ln(1-x)\le x/(1-x)$ for $0\le x<1$, the expression in parentheses is at most
\[
 -\frac14+\frac{49}{2304}+\frac{49}{1250}<0.
\]
This proves \eqref{eq:two-parent} in the interior. For the boundary cases, approximate the parent mean by interior means and add an arbitrarily small independent fair-bit noise to the child. Conditional entropy and all terms in \eqref{eq:two-parent} are continuous in the resulting finite joint laws, so the inequality passes to the limit.

Finally, put $a=9/25$. Then $w=a/(1-a)$ and $\tau-1/4=a/2$. Consequently,
\[
 \kappa+\sigma\tau
 =\frac12+\frac12\bigl(\log_2(1+w)-a\log_2 w\bigr)
 =\frac{1+h(a)}2.
\]
The first line of \eqref{eq:scalar-bounds} proves \eqref{eq:theta-identity}.
\end{proof}

\subsection{Sampling parents without replacement}

\begin{lemma}\label{lem:without-replacement}
Consider an array of $L\ge4$ bits with mean $q$. Sample an ordered pair of distinct indices uniformly, and let $X,Y$ be the sampled bits. For an arbitrary conditional law of a binary child $Z$ given $(X,Y)$, put $v=\E Z$ and let $d$ be as in Lemma~\ref{lem:two-parent}. Then
\begin{equation}\label{eq:without-replacement}
 \HH(Z\mid X,Y)\le
 \kappa+\sigma d+\frac{h(v)-h(q)}2+\eta(L),
\end{equation}
where
\begin{equation}\label{eq:eta}
 \eta(L)=\frac{1+\sigma}{L}+\frac12h(1/L).
\end{equation}
\end{lemma}

\begin{proof}
Compare the parent law with the law obtained by sampling two indices independently with replacement. These laws have total variation distance at most $1/L$. One coupling samples both indices independently and, only when they coincide, resamples the second index uniformly from the other $L-1$ choices. The resulting distinct pair is uniform.

Use the same child kernel under the two parent laws, assigning it arbitrarily at parent types of zero mass. Conditional entropy is the expectation of a function in $[0,1]$, so it changes by at most $1/L$. The disagreement cost changes by at most $1/L$, and so does the child mean. The parent mean remains $q$. By \eqref{eq:continuity}, the two child entropies differ by at most $h(1/L)$, since $h$ is increasing on $[0,1/2]$. Applying Lemma~\ref{lem:two-parent} to the with-replacement law gives \eqref{eq:without-replacement}.
\end{proof}

\section{A fixed locally sparse graph}

Fix
\begin{equation}\label{eq:fixed-parameters}
 L_0=10^7,\qquad \delta=10^{-5},\qquad t=20000.
\end{equation}
Let $V_0$ be a set of $L_0$ vertices. Having constructed $V_{i-1}$, let $V_i$ be a disjoint set indexed by the unordered pairs of distinct vertices of $V_{i-1}$. Join each vertex of $V_i$ to the two vertices indexing it. Do this for $1\le i\le t$, and include exactly these edges. Denote the resulting graph by $G$. Write $L_i=\abs{V_i}$, so that
\[
 L_i=\binom{L_{i-1}}2.
\]
Every $L_i$ is finite, and $L_i\ge L_0$. In particular, $G$ is fixed independently of every parameter introduced below.

\begin{lemma}\label{lem:local-sparsity}
The graph $G$ is bipartite and $2$-degenerate, and satisfies \eqref{eq:local-condition}.
\end{lemma}

\begin{proof}
Even and odd layers give a bipartition. Order the vertices by increasing layer, with any order inside a layer. Every vertex has at most two neighbors earlier in this order, which proves $2$-degeneracy. Restrict the order to a set $W$ of $s\ge2$ vertices. Its first vertex has no earlier neighbor, its second has at most one, and each remaining vertex has at most two. Counting edges by their later endpoint gives
\[
 e(G[W])\le0+1+2(s-2)=2s-3.
\]
\end{proof}

The next elementary bounds ensure a uniform entropy margin at every layer.

\begin{lemma}\label{lem:fixed-margins}
For every integer $L\ge L_0$, with $M=\binom L2$,
\begin{equation}\label{eq:margins}
 \eta(L)<\delta,\qquad
 L+3\log_2(M+1)-\delta M<-1.
\end{equation}
Also,
\begin{equation}\label{eq:depth-margin}
 2t(\beta-\theta-2\delta)>1.
\end{equation}
\end{lemma}

\begin{proof}
We have $0<\sigma<1$. For $0<x<1$, the inequality $-(1-x)\ln(1-x)\le x$ gives
\[
 h(x)\le x\log_2(1/x)+\frac{x}{\ln2}.
\]
Using $\log_2 L_0<24$ and $1/\ln2<2$ yields
\[
 \eta(L_0)<\frac2{10^7}+\frac{26}{2\cdot10^7}
 =\frac{15}{10^7}<\delta.
\]
Both terms in \eqref{eq:eta} decrease for $L\ge L_0$, so the first bound follows. Next, $M+1\le L^2$ and $\log_2L\le L$. Furthermore,
\[
 \delta M=\frac{L(L-1)}{200000}\ge49L.
\]
Thus the second expression in \eqref{eq:margins} is at most $L+6L-49L=-42L<-1$. Finally, \eqref{eq:theta-identity} gives
\[
 \beta-\theta-2\delta>.9714-.97135-.00002=.00003.
\]
Multiplying by $2t=40000$ proves \eqref{eq:depth-margin}.
\end{proof}

\section{Excluding the fixed graph in a sampled cube}

From now on, $N$ tends to infinity through positive multiples of $40000$. Let
\begin{equation}\label{eq:cube-parameters}
 U_N=\{0,1\}^N,\qquad Q=2^N,\qquad
 k=\tau N,\qquad p=2^{-\beta N}.
\end{equation}
Then $k$ is an even integer. The graph $B_N$ on $U_N$ joins two distinct words when their Hamming distance is at most $k$. Each word is retained independently with probability $p$ to form a random set $S$. We shall prove that $B_N[S]$ is $G$-free with probability tending to one.

\subsection{Counting child arrays}

Consider a single transition with $L$ parents and $M=\binom L2$ children. Fix an ordered parent array $u=(u_1,\ldots,u_L)\in U_N^L$. A child array $z\in U_N^M$ assigns a word to each unordered pair of parent indices.

For coordinate $j$, choose a child index uniformly and orient its two parents by a fair coin. Denote the parent bits by $X_j,Y_j$ and the child bit by $Z_j$. Define
\begin{equation}\label{eq:profile-entropy}
 E(u,z)=\frac1N\sum_{j=1}^N\HH(Z_j\mid X_j,Y_j).
\end{equation}

\begin{lemma}\label{lem:array-count}
For a fixed parent array $u$, the number of child arrays satisfying $E(u,z)\le\beta-\delta$ is at most
\begin{equation}\label{eq:array-count}
 (M+1)^{3N}2^{NM(\beta-\delta)}.
\end{equation}
\end{lemma}

\begin{proof}
At each coordinate, split the child indices into the three groups corresponding to unordered parent types $00$, $01$, and $11$. Let their sizes be $M_{00},M_{01},M_{11}$, and let $a_{00},a_{01},a_{11}$ count the child bits equal to one in these groups. The three counts have at most $(M+1)^3$ possibilities. For a specified triple, the number of child-bit assignments is
\[
 \prod_{g\in\{00,01,11\}}\binom{M_g}{a_g}
 \le 2^{\sum_g M_g h(a_g/M_g)},
\]
by \eqref{eq:binomial-entropy}, with an empty group contributing zero to the exponent. The fair orientation of a mixed parent pair gives the same conditional child distribution for parent types $01$ and $10$. Hence the exponent is exactly $M\HH(Z_j\mid X_j,Y_j)$.

Specify such a triple at every coordinate. There are at most $(M+1)^{3N}$ resulting profiles, and for each profile the total number of arrays is at most $2^{NM E(u,z)}$. Summing over the profiles with $E(u,z)\le\beta-\delta$ proves the lemma.
\end{proof}

\subsection{The entropy potential across layers}

\begin{proposition}\label{prop:red-free}
For the fixed graph $G$ of Section~3,
\begin{equation}\label{eq:red-failure}
 \Prob\bigl(G\subseteq B_N[S]\bigr)\le t2^{-N}.
\end{equation}
\end{proposition}

\begin{proof}
For a transition with $L$ parents and $M$ children, call a child array distinct if its $M$ words are pairwise distinct. Every particular distinct child array is contained in $S$ with probability $p^M$. There are $Q^L$ possible parent arrays. Lemmas~\ref{lem:array-count} and \ref{lem:fixed-margins}, followed by the union bound, show that the probability of a parent array and a distinct surviving child array with $E(u,z)\le\beta-\delta$ is at most
\begin{align*}
 Q^L(M+1)^{3N}2^{NM(\beta-\delta)}p^M
 &=2^{N[L+3\log_2(M+1)-\delta M]}\\
 &\le2^{-N}.
\end{align*}
This calculation counts all parent arrays, regardless of whether their words survive. Taking a union bound over the $t$ transitions gives an event of probability at least $1-t2^{-N}$ on which every distinct surviving child array at every transition has entropy greater than $\beta-\delta$.

Suppose this event holds and there is an injective copy of $G$ in $B_N[S]$. For coordinate $j$, let $q_{i,j}$ be the proportion of ones among the images of the vertices in layer $V_i$, and put
\[
 \Phi_i=\frac1N\sum_{j=1}^N h(q_{i,j}).
\]
We have $0\le\Phi_i\le1$. For the transition from $V_{i-1}$ to $V_i$, let $E_i$ be the quantity \eqref{eq:profile-entropy}. The child images are distinct and survive, so $E_i>\beta-\delta$.

Let $d_{i,j}$ be the average of the two parent-child disagreement probabilities at coordinate $j$. Every parent-child edge of the copy has Hamming distance at most $k$. Averaging over all $2L_i$ parent-child edges therefore gives
\[
 \frac1N\sum_{j=1}^N d_{i,j}\le\frac{k}{N}=\tau.
\]
At each coordinate the oriented parent pair is a uniformly sampled pair of distinct indices from the parent layer. Apply Lemma~\ref{lem:without-replacement} and average over coordinates. By \eqref{eq:theta-identity},
\[
 E_i\le\theta+\frac{\Phi_i-\Phi_{i-1}}2+\eta(L_{i-1})
 <\theta+\frac{\Phi_i-\Phi_{i-1}}2+\delta.
\]
Thus
\[
 \Phi_i-\Phi_{i-1}>2(\beta-\theta-2\delta).
\]
Summing for $1\le i\le t$ and using \eqref{eq:depth-margin} yields $\Phi_t-\Phi_0>1$, contradicting $0\le\Phi_0,\Phi_t\le1$. Therefore the event under consideration is $G$-free, which proves \eqref{eq:red-failure}.
\end{proof}

\section{The spectrum of the Hamming ball}

We next control independent sets of $B_N$. Let $T$ be the real symmetric matrix indexed by $U_N$ with
\[
 T(x,y)=\one_{\dist(x,y)\le k}.
\]
Its diagonal entries are one and its constant row sum is
\begin{equation}\label{eq:degree}
 D=\sum_{j=0}^k\binom Nj.
\end{equation}
By \eqref{eq:binomial-entropy} and the monotonicity of $h$ on $[0,1/2]$,
\begin{equation}\label{eq:D-bounds}
 \frac{2^{Nh(\tau)}}{N+1}\le D\le(N+1)2^{Nh(\tau)}.
\end{equation}
In particular $D=2^{(h(\tau)+o(1))N}$.

\subsection{Characters and Krawtchouk polynomials}

Identify $U_N$ with $\F_2^N$ and write $\abs x$ for the Hamming weight of a word $x$. For $b\in\F_2^N$, define
\[
 \chi_b(x)=(-1)^{b\cdot x}.
\]
The characters are orthogonal under summation over $U_N$, each with squared norm $Q$. As $T(x,y)$ depends only on $x+y$, each character is an eigenvector. Define $K_j^{(n)}(v)$ by
\begin{equation}\label{eq:K-generating}
 (1-z)^v(1+z)^{n-v}=\sum_{j=0}^n K_j^{(n)}(v)z^j.
\end{equation}
Equivalently,
\[
 K_j^{(n)}(v)=\sum_{a=0}^j(-1)^a\binom va\binom{n-v}{j-a}.
\]
Using the polynomial definition of the binomial coefficients extends $K_j^{(n)}$ to a polynomial in the real variable $v$.
If $\abs{b}=i$, summing $\chi_b$ over words of a given weight shows that its eigenvalue is
\[
 \lambda_i=\sum_{j=0}^k K_j^{(N)}(i).
\]
We have $\lambda_0=D$. For $1\le i\le N$, coefficient extraction gives
\begin{equation}\label{eq:partial-sum}
 \lambda_i=[z^k]\frac{(1-z)^i(1+z)^{N-i}}{1-z}
 =K_k^{(N-1)}(i-1).
\end{equation}

\begin{lemma}\label{lem:root-location}
Let $1\le k\le(n+1)/2$ be even. If $K_k^{(n)}(v)<0$ for an integer $0\le v\le n$, then
\begin{equation}\label{eq:root-location}
 \abs{n-2v}<2\sqrt{k(n-k+1)}.
\end{equation}
\end{lemma}

\begin{proof}
Put $x=n-2v$. Differentiating \eqref{eq:K-generating} and comparing coefficients gives
\[
 (j+1)K_{j+1}=xK_j-(n-j+1)K_{j-1}.
\]
Thus the monic polynomials $P_j(x)=j!K_j^{(n)}((n-x)/2)$ satisfy
\begin{equation}\label{eq:monic-recurrence}
 P_0=1,\qquad P_1=x,\qquad
 P_{j+1}=xP_j-j(n-j+1)P_{j-1}.
\end{equation}
The determinant recurrence identifies $P_k$ with the characteristic polynomial of the real symmetric $k\times k$ tridiagonal matrix having zero diagonal and off-diagonal entries
\[
 \sqrt{j(n-j+1)},\qquad 1\le j\le k-1.
\]
All roots are real. Each eigenvalue has absolute value at most the maximal absolute row sum, which is at most $2\sqrt{k(n-k+1)}$. Here $j(n-j+1)$ is increasing for $0\le j\le(n+1)/2$.

The monic polynomial $P_k$ is positive to the right of its largest root. Recurrence \eqref{eq:monic-recurrence} also gives $P_k(-x)=(-1)^kP_k(x)$. Since $k$ is even, $P_k$ is positive on both exterior intervals beyond the indicated root bound, and is nonnegative at the endpoints. A negative value must satisfy \eqref{eq:root-location}.
\end{proof}

\begin{lemma}\label{lem:K-parseval}
For $0\le k\le n$,
\begin{equation}\label{eq:K-parseval}
 \sum_{v=0}^n\binom nv K_k^{(n)}(v)^2=2^n\binom nk.
\end{equation}
\end{lemma}

\begin{proof}
On $\F_2^n$, let $f(x)=\one_{\abs{x}=k}$ and define the unnormalized Fourier coefficient $\widehat f(b)=\sum_x f(x)\chi_b(x)$. If $\abs b=v$, then \eqref{eq:K-generating} gives $\widehat f(b)=K_k^{(n)}(v)$. Orthogonality gives
\[
 \sum_b\widehat f(b)^2
 =\sum_{x,y}f(x)f(y)\sum_b\chi_b(x+y)
 =2^n\sum_x f(x)^2
 =2^n\binom nk.
\]
Grouping $b$ by its weight proves \eqref{eq:K-parseval}.
\end{proof}

\subsection{A bound on the negative eigenvalues}

\begin{proposition}\label{prop:negative-spectrum}
For all sufficiently large admissible $N$, every eigenvalue of $T$ is at least $-\Lambda$, where
\begin{equation}\label{eq:spectral-decay}
 \frac{\Lambda}{D}\le2^{-.0152N}.
\end{equation}
\end{proposition}

\begin{proof}
Set $n=N-1$. Since $k=\tau N$ is even and $k\le(n+1)/2$, Lemma~\ref{lem:root-location} applies. If the eigenvalue in \eqref{eq:partial-sum} is negative, with $v=i-1$, then
\[
 \frac12-\frac{\sqrt{k(n-k+1)}}n
 <\frac vn<
 \frac12+\frac{\sqrt{k(n-k+1)}}n.
\]
Now
\[
 \frac{\sqrt{k(n-k+1)}}n
 =\frac{N}{N-1}\sqrt{\tau(1-\tau)}.
\]
The exact inequality $.4951^2=.24512401>.2451=\tau(1-\tau)$ implies
\[
 \frac12-\sqrt{\tau(1-\tau)}>.0049.
\]
Consequently every negative eigenvalue has $v/n\in[.0049,.9951]$ for all sufficiently large $N$.

By Lemmas~\ref{lem:elementary-entropy} and \ref{lem:K-parseval},
\begin{align*}
 \abs{K_k^{(n)}(v)}^2
 &\le \frac{2^n\binom nk}{\binom nv}\\
 &\le(n+1)\,2^{n[1+h(k/n)-h(v/n)]}\\
 &\le(n+1)\,2^{n[1+h(k/n)-h(.0049)]}.
\end{align*}
In the last step we use the symmetry and monotonicity of binary entropy. Since $k/n\to\tau$, this bounds the magnitude of every negative eigenvalue by
\[
 2^{\left(\frac{1+h(\tau)-h(.0049)}2+o(1)\right)N}.
\]
After division by the lower bound in \eqref{eq:D-bounds}, the exponent becomes
\[
 -\frac{h(\tau)+h(.0049)-1}{2}+o(1).
\]
The estimates \eqref{eq:scalar-bounds} give
\[
 \frac{h(\tau)+h(.0049)-1}{2}>.015225>.0152.
\]
Taking $\Lambda$ to be the larger of zero and the magnitude of the most negative eigenvalue proves \eqref{eq:spectral-decay}.
\end{proof}

\section{Independent-set containers and random thinning}

Put
\begin{equation}\label{eq:container-parameters}
 M_N=\ceil{2^{.985N}},\qquad
 b_N=2^{.01425N}.
\end{equation}
For admissible $N=40000m$, we have $b_N=2^{570m}$.

\subsection{Supersaturation}

\begin{lemma}\label{lem:supersaturation}
For all sufficiently large admissible $N$, every $W\subseteq U_N$ with $\abs W\ge M_N$ contains a vertex whose degree in $B_N[W]$ is at least
\begin{equation}\label{eq:large-degree}
 \frac{D\abs W}{4Q}.
\end{equation}
\end{lemma}

\begin{proof}
Decompose $\one_W=(\abs W/Q)\one+f$, where $f$ is orthogonal to the constant vector. Proposition~\ref{prop:negative-spectrum} and $\norm f^2=\abs W-\abs W^2/Q\le\abs W$ give
\[
 \one_W^{\mathsf T}T\one_W
 \ge \frac{D\abs W^2}{Q}-\Lambda\abs W.
\]
Since $\abs W/Q\ge2^{-.015N}$ and $\Lambda/D\le2^{-.0152N}$, for large $N$ the second term is at most half the first. Hence
\[
 \one_W^{\mathsf T}T\one_W\ge\frac{D\abs W^2}{2Q}.
\]
There is one diagonal entry for each vertex, so
\[
 \frac{2e(B_N[W])}{\abs W}
 =\frac{\one_W^{\mathsf T}T\one_W-\abs W}{\abs W}
 \ge\frac{D\abs W}{2Q}-1.
\]
By \eqref{eq:D-bounds}, $DM_N/Q\to\infty$. The last expression is therefore at least $D\abs W/(4Q)$, uniformly for $\abs W\ge M_N$. Some vertex has at least the average degree.
\end{proof}

\subsection{A deterministic encoding}

Fix a total order of $U_N$ to resolve ties between vertices of the same degree. Given an independent set $I$ of $B_N$, start with $W=U_N$ and an empty set $F$. While $\abs W\ge M_N$, take a vertex $x$ of maximum degree in $B_N[W]$ and do the following.
\begin{enumerate}
 \item If $x\notin I$, delete $x$ from $W$.
 \item If $x\in I$, add $x$ to $F$ and delete $x$ and all its current neighbors from $W$.
\end{enumerate}
Write $W_F$ for the set remaining when the procedure stops. Independence implies
\begin{equation}\label{eq:container-cover}
 I\subseteq F\cup W_F,\qquad \abs{W_F}<M_N.
\end{equation}

\begin{lemma}\label{lem:fingerprints}
For all sufficiently large admissible $N$, the set $W_F$ is determined by the unordered set $F$. Moreover, every set $F$ produced by the procedure has size at most
\begin{equation}\label{eq:fingerprint-size}
 L_N=\ceil{1+\frac{4Q}{D}\ln Q}.
\end{equation}
The number of possible pairs $(F,W_F)$ is at most $(Q+1)^{L_N+1}$, and
\begin{equation}\label{eq:fingerprint-negligible}
 L_N=o(b_N),\qquad (L_N+1)\ln(Q+1)=o(b_N).
\end{equation}
\end{lemma}

\begin{proof}
Replay the procedure using membership in $F$ to choose between its two deletion rules. At each stage of the original run, the selected vertex belongs to $I$ exactly when it belongs to $F$. Indeed, every selected vertex of $I$ is recorded, and all recorded vertices lie in $I$. An induction shows that the replay selects the same vertices and deletes the same sets. This proves the determination assertion.

Put $a=D/(4Q)$, so $0<a\le1/4$. Whenever a vertex is recorded, Lemma~\ref{lem:supersaturation} shows that the current set loses at least the fraction $a$ of its vertices. Other steps also decrease its size. If $s\ge1$ vertices are recorded, the set immediately before the last such step has size at least $M_N$ and at most $Q(1-a)^{s-1}$. Therefore
\[
 s-1\le\frac{\ln(Q/M_N)}{-\ln(1-a)}
 \le\frac{\ln Q}{a}=\frac{4Q}{D}\ln Q,
\]
which gives \eqref{eq:fingerprint-size}. The number of subsets of $U_N$ of size at most $L_N$ is at most $(Q+1)^{L_N+1}$. Each such recorded set determines at most one pair $(F,W_F)$.

Finally, \eqref{eq:D-bounds} yields
\[
 L_N=O\left(N^2 2^{(1-h(\tau))N}\right).
\]
As $1-h(\tau)<.01419<.01425$, both assertions of \eqref{eq:fingerprint-negligible} follow. A polynomial factor is dominated by the exponential gap between these two constants.
\end{proof}

\subsection{Thinning the containers}

\begin{proposition}\label{prop:blue-free}
For the independently retained set $S$ with retention probability $p=2^{-\beta N}$,
\begin{equation}\label{eq:blue-failure}
 \Prob\bigl(\alpha(B_N[S])\ge b_N\bigr)=o(1).
\end{equation}
\end{proposition}

\begin{proof}
Every $W_F$ in Lemma~\ref{lem:fingerprints} is a fixed subset of the ambient cube. Let $X_F=\abs{W_F\cap S}$. Independence of the retention variables gives the exact exponential moment
\begin{equation}\label{eq:exponential-moment}
 \E 2^{X_F}=(1+p)^{\abs{W_F}}\le\exp(pM_N).
\end{equation}
The parameters satisfy
\[
 pM_N\le2^{(.985-\beta)N}+1
 =2^{.0136N}+1=o(b_N).
\]
For large $N$, \eqref{eq:fingerprint-negligible} gives $L_N\le b_N/2$, and the preceding bound gives $pM_N\le(\ln2)b_N/4$. Markov's inequality applied to \eqref{eq:exponential-moment} then yields
\begin{align*}
 \Prob(X_F\ge b_N-L_N)
 &\le \exp\bigl(pM_N-(\ln2)(b_N-L_N)\bigr)\\
 &\le\exp\bigl(- (\ln2)b_N/4\bigr).
\end{align*}
If there is an independent set $I\subseteq S$ of size at least $b_N$, its encoding satisfies \eqref{eq:container-cover}. Hence $X_F\ge b_N-\abs F\ge b_N-L_N$. Taking the union bound over all possible pairs $(F,W_F)$ gives
\[
 \Prob\bigl(\alpha(B_N[S])\ge b_N\bigr)
 \le\exp\left((L_N+1)\ln(Q+1)-\frac{\ln2}{4}b_N\right)=o(1),
\]
by \eqref{eq:fingerprint-negligible}.
\end{proof}

\section{Completion of the Ramsey argument}

\begin{proof}[Proof of Theorem~\ref{thm:main}]
Lemma~\ref{lem:local-sparsity} gives the required structural properties of the fixed graph $G$. It remains to prove the Ramsey bound.

The order of the random sample has mean $pQ=2^{(1-\beta)N}$ and variance at most $pQ$. Chebyshev's inequality gives
\begin{equation}\label{eq:sample-size}
 \Prob\left(\abs S<\frac{pQ}{2}\right)
 \le\frac4{pQ}=o(1).
\end{equation}
Together with Proposition~\ref{prop:red-free} and Proposition~\ref{prop:blue-free}, this shows that, for every sufficiently large admissible $N$, there exists a set $S\subseteq U_N$ satisfying all three properties
\begin{equation}\label{eq:three-properties}
 \abs S\ge\frac12 2^{(1-\beta)N},\qquad
 G\nsubseteq B_N[S],\qquad
 \alpha(B_N[S])<b_N.
\end{equation}
Indeed, the sum of the three failure probabilities tends to zero.

Color each edge of the complete graph on $S$ red if its endpoints are adjacent in $B_N$, and blue otherwise. The second property in \eqref{eq:three-properties} excludes a red $G$. A blue clique is an independent set in $B_N[S]$, so the third property excludes a blue $K_{b_N}$. By the definition of the Ramsey number,
\[
 R(G,K_{b_N})>\abs S\ge\frac12 2^{(1-\beta)N}.
\]
Taking $N=40000m$ gives
\[
 b_N=2^{570m},\qquad (1-\beta)N=1144m,
\]
and proves \eqref{eq:main-bound}. For these cliques,
\[
 e\bigl(K_{2^{570m}}\bigr)
 =\frac{2^{570m}(2^{570m}-1)}2<2^{1140m-1}.
\]
Division proves \eqref{eq:main-ratio}. The cliques have no isolated vertices. Their ratios tend to infinity while $G$ remains fixed, which disproves \eqref{eq:size-linear} for a graph satisfying the full local hypothesis \eqref{eq:local-condition}.
\end{proof}

\appendix
\section{Exact scalar inequalities}\label{app:scalars}

This appendix gives rational certificates for \eqref{eq:scalar-bounds}. For $1\le y\le2$, set
\[
 s(y)=2\sum_{j=0}^{15}\frac{1}{2j+1}
       \left(\frac{y-1}{y+1}\right)^{2j+1},
 \qquad \rho=\frac{3}{44\cdot3^{33}}.
\]
The expansion $\ln y=2\operatorname{artanh}((y-1)/(y+1))$ has nonnegative terms. Since $(y-1)/(y+1)\le1/3$, its omitted tail satisfies
\[
 0\le\ln y-s(y)
 \le\frac{2(1/3)^{33}}{33(1-1/9)}=\rho<10^{-16}.
\]
Thus
\begin{equation}\label{eq:log-enclosure}
 s(y)\le\ln y\le s(y)+\rho.
\end{equation}
Let $\ell=s(2)$. For a rational $x=2^e y\in(0,1)$ with $1\le y\le2$ and $e<0$, define
\[
 L_x=s(y)+e(\ell+\rho),\qquad
 U_x=s(y)+\rho+e\ell.
\]
Equation \eqref{eq:log-enclosure} implies $L_x\le\ln x\le U_x$. Define the rational numbers
\[
 A_x=-xU_x-(1-x)U_{1-x},\qquad
 B_x=-xL_x-(1-x)L_{1-x}.
\]
For the three arguments below, $A_x>0$, and therefore
\begin{equation}\label{eq:entropy-enclosure}
 \frac{A_x}{\ell+\rho}\le h(x)\le\frac{B_x}{\ell}.
\end{equation}
The required power-of-two reductions are
\[
\begin{array}{c|cc|cc}
 x&e&y&e'&y'\\ \hline
 9/25&-2&36/25&-1&32/25\\
 43/100&-2&43/25&-1&57/50\\
 49/10000&-8&784/625&-1&9951/5000
\end{array}
\]
where $x=2^e y$ and $1-x=2^{e'}y'$. To make the final comparisons explicit, take
\[
\begin{array}{c|cc}
 x&a_x&b_x\\ \hline
 9/25&94268/100000&94270/100000\\
 43/100&98581/100000&98582/100000\\
 49/10000&4464/100000&4465/100000.
\end{array}
\]
Substitution in the finite rational sums gives, in all three rows,
\begin{equation}\label{eq:rational-certificates}
 A_x-a_x(\ell+\rho)>\frac1{10^7},\qquad
 b_x\ell-B_x>\frac1{10^7}.
\end{equation}
These inequalities involve rational numbers alone and follow by clearing their positive denominators. They imply $A_x>0$ and, together with \eqref{eq:entropy-enclosure}, prove $a_x<h(x)<b_x$. This is exactly \eqref{eq:scalar-bounds}.

For completeness, the other scalar estimates used in the proof follow from
\[
 \frac{49}{2304}+\frac{49}{1250}<\frac14,
 \qquad \left(\frac{4951}{10000}\right)^2>\frac{43\cdot57}{10000},
\]
and the direct consequences
\[
 \frac{h(.43)+h(.0049)-1}{2}>.015225,
 \qquad 1-h(.43)<.01419.
\]
All decimal constants in the proof are terminating decimals and hence denote exact rational numbers.

\section*{Acknowledgment}
The authors acknowledge assistance from OpenAI's ChatGPT in developing and checking the argument and in preparing the exposition.

\begin{thebibliography}{9}

\bibitem{BTW}
J. Balogh, A. Treglown, and A. Z. Wagner,
\emph{Applications of graph containers in the Boolean lattice},
Random Structures Algorithms \textbf{49} (2016), no. 4, 845--872.
\href{https://doi.org/10.1002/rsa.20666}{doi:10.1002/rsa.20666}.

\bibitem{Bloom}
T. F. Bloom,
\emph{Erd\H{o}s problems, Problem 566},
\url{https://www.erdosproblems.com/566}, accessed November 11, 2026.

\bibitem{EFRS}
P. Erd\H{o}s, R. J. Faudree, C. C. Rousseau, and R. H. Schelp,
\emph{Ramsey size linear graphs},
Combin. Probab. Comput. \textbf{2} (1993), no. 4, 389--399.
\href{https://doi.org/10.1017/S096354830000078X}{doi:10.1017/S096354830000078X}.

\bibitem{OpenAI}
OpenAI,
\emph{Ten advances in mathematics and theoretical computer science},
preprint, updated August 6, 2026, Chapter 10, pp. 236--248.
\url{https://cdn.openai.com/pdf/ten-proofs-oai.pdf}.

\end{thebibliography}
\end{document}
